\documentclass[12pt]{article}

% --- 1. CODIFICACIÓN E IDIOMA ---
\usepackage[utf8]{inputenc}
\usepackage[spanish,es-nodecimaldot]{babel} % es-nodecimaldot evita conflictos con siunitx

% --- 2. GEOMETRÍA Y DISEÑO ---
\usepackage{geometry}
\geometry{a4paper, margin=1in}
\usepackage{pdflscape} % Para páginas horizontales si las necesitas

% --- 3. COLORES Y TABLAS (Orden Crítico) ---
% xcolor debe cargarse ANTES que otros paquetes que usen color
\usepackage[table,xcdraw]{xcolor}
\usepackage{booktabs}   % Tablas profesionales
\usepackage{multirow}
\usepackage{array}
\usepackage{tabularx}
\usepackage{longtable}
\usepackage{float}      % Para usar [H]

% --- 4. MATEMÁTICAS Y CIENCIA ---
\usepackage{amsmath, amssymb, amsthm}
\usepackage{siunitx}    % Para unidades (kg, kWh, etc.)

% --- 5. GRÁFICOS Y TIKZ ---
\usepackage{graphicx}
\usepackage{caption}
\usepackage{subcaption}
\usepackage{tikz}
% Cargamos todas las librerías necesarias de una sola vez
\usetikzlibrary{
  shapes,
  shapes.geometric,
  arrows,
  positioning,
  babel,
  calc,
  decorations.pathmorphing,
  patterns,
  shadows
}


% --- 6. CÓDIGO FUENTE ---
\usepackage{listings}

% --- 7. HIPERVÍNCULOS (Cargar al final) ---

\usepackage{hyperref}
\hypersetup{
  colorlinks=true,
  linkcolor=blue,
  filecolor=magenta,
  urlcolor=cyan,
  pdfpagemode=FullScreen
}
% --- DEFINICIÓN DE COLORES PERSONALIZADOS ---
\definecolor{bluebg}{RGB}{230,240,255}
\definecolor{greenbg}{RGB}{230,255,230}
\definecolor{graybg}{RGB}{240,240,240}
\definecolor{concepto}{RGB}{0, 102, 204}
\definecolor{ejemplo}{RGB}{0, 153, 76}
\definecolor{importante}{RGB}{204, 0, 0}
\definecolor{sustrato_aerobio}{HTML}{D7B56D}
\definecolor{sustrato_anoxico}{HTML}{5D4C3D}
\definecolor{larva_color}{HTML}{FCEEB5}
\definecolor{gas_bueno}{HTML}{4682B4}
\definecolor{gas_neutro}{HTML}{A9A9A9}
\definecolor{gas_malo_ch4}{HTML}{FF4500}
\definecolor{gas_malo_n2o}{HTML}{8B0000}

% --- CONFIGURACIÓN DE LISTINGS (CÓDIGO) ---
\lstset{
    language=Python,
    basicstyle=\ttfamily\small,
    keywordstyle=\color{blue},
    commentstyle=\color{ejemplo}, % Usamos tu verde definido
    stringstyle=\color{importante}, % Usamos tu rojo definido
    showstringspaces=false,
    breaklines=true,
    frame=single,
    backgroundcolor=\color{graybg},
    rulecolor=\color{gray},
    tabsize=4,
    numbers=left,
    numberstyle=\tiny\color{gray}
}

% --- METADATOS DEL DOCUMENTO ---
\title{Análisis Ambiental de Proceso de Bioconversión de Mosca Soldado Negra \\ \large Sensor Virtual con PGNN para Evaluación Ambiental}
\author{Propuesta Técnica Detallada}
\date{\today}
\begin{document}

\maketitle

\begin{abstract}
\noindent Esta propuesta detalla la implementacion de un sensor virtual basado en Physics-Guided Neural Networks (PGNN) para la evaluacion ambiental de procesos de bioconversion larval utilizando \textit{Hermetia illucens} o \textit{Tenebrio molitor}. El documento presenta una metodologia completa que integra fundamentos bioquimicos, formulaciones matematicas y protocolos experimentales especificos para la cuantificacion de gases de efecto invernadero (GEI) y otros indicadores ambientales clave.
\end{abstract}

\tableofcontents

\newpage

\section{Fase 1: Definicion del Indicador Ambiental (KPI) -- Detalle Tecnico}

\subsection{Marco de evaluacion: ACV simplificado con enfoque en GEI}

\begin{table}[ht]
\centering
\caption{Componentes del sistema y flujos a cuantificar para el proceso de \textit{Hermetia illucens}}
\label{tab:componentes_sistema}
\small
\begin{tabularx}{\textwidth}{@{} l l l X @{}}
\toprule
\textbf{Componente} & \textbf{Flujo} & \textbf{Unidad} & \textbf{Método de cálculo} \\ \midrule
\rowcolor[gray]{0.95} \multicolumn{4}{l}{\textit{Entradas}} \\
Directas & Residuo orgánico & kg MS/día & Balanza + secado (\SI{105}{\celsius}, \SI{24}{\hour}) \\ \midrule
\rowcolor[gray]{0.95} \multicolumn{4}{l}{\textit{Salidas Directas}} \\
Emisiones & $CO_{2}$ emitido & kg/día & Sensor NDIR + integración temporal \\
& $CH_{4}$ emitido & kg/día & Sensor electroquímico o GC \\
& $N_{2}O$ emitido* & kg/día & Inferido vía PGNN (relación N residual) \\
Productos & Biomasa larval & kg MS/día & Cosecha + secado \\
& Frass (residuo) & kg MS/día & Pesaje post-proceso \\ \midrule
\rowcolor[gray]{0.95} \multicolumn{4}{l}{\textit{Entradas Indirectas}} \\
Recursos & Energía eléctrica & kWh/día & Medidor de consumo (vatios) \\
& Agua & L/día & Contador de flujo \\ \bottomrule
\end{tabularx}
\end{table}

\noindent\textbf{Nota:} N\textsubscript{2}O es critico (265$\times$ GWP de CO\textsubscript{2}) pero dificil de medir; debe inferirse via relacion con nitrogeno amoniacal residual y pH \cite{ref59}.


\section{Explicación Detallada del Marco de Evaluación ACV Simplificado}

\subsection{¿Qué es el Análisis de Ciclo de Vida (ACV)?}

El \textbf{Análisis de Ciclo de Vida (ACV)}, también conocido como \textit{Life Cycle Assessment (LCA)} en inglés, es una metodología sistemática para evaluar los impactos ambientales asociados con todas las etapas de la vida de un producto, proceso o servicio. En el contexto de la bioconversión larval, el ACV nos permite cuantificar y comparar la huella ambiental del proceso frente a alternativas tradicionales de manejo de residuos.

\begin{center}
\begin{tikzpicture}[node distance=0.5cm]
\node[draw, rectangle, fill=blue!10, rounded corners, minimum width=12cm, minimum height=1.5cm, align=center, text width=11cm] (def) {
\textbf{ACV Simplificado para Bioconversión} \\
Evalúa los flujos de materia y energía específicamente relacionados con las emisiones de gases de efecto invernadero (GEI) dentro de los límites de la planta (Gate-to-Gate).
};
\end{tikzpicture}
\end{center}

\subsection{Objetivo del Marco de Evaluación}

El objetivo principal de este marco es \textbf{cuantificar las emisiones netas de gases de efecto invernadero} generadas durante el proceso de bioconversión de residuos orgánicos mediante larvas de \textit{Hermetia illucens}, permitiendo:

\begin{itemize}
    \item \textcolor{concepto}{\textbf{Comparación objetiva}} con procesos alternativos (compostaje, digestión anaeróbica).
    \item \textcolor{concepto}{\textbf{Identificación de puntos críticos}} ("Hotspots") donde se generan mayores emisiones.
    \item \textcolor{concepto}{\textbf{Optimización del proceso}} mediante control de variables operativas (aireación, humedad).
    \item \textcolor{concepto}{\textbf{Certificación ambiental}} alineada con estándares (GHG Protocol, ISO 14067).
\end{itemize}

\section{Explicación Detallada de la Tabla de Componentes}

La tabla organizada en tres grandes categorías refleja el \textbf{balance de masa y energía} del sistema:

\begin{table}[H]
\centering
\small
\caption{Componentes del sistema y flujos a cuantificar}
\label{tab:componentes_sistema}
\begin{tabularx}{\textwidth}{@{} l l l X @{}}
\toprule
\textbf{Componente} & \textbf{Flujo} & \textbf{Unidad} & \textbf{Método de cálculo} \\
\midrule
\multicolumn{4}{l}{\textit{\textbf{Entradas}}} \\
Materia Prima & Residuo orgánico & kg MS/día & Balanza + determinación MS (\SI{105}{\celsius}, \SI{24}{\hour}) \\
\midrule
\multicolumn{4}{l}{\textit{\textbf{Salidas Directas}}} \\
Emisiones & $CO_{2}$ emitido & kg/día & Sensor NDIR + integración temporal \\
          & $CH_{4}$ emitido & kg/día & Sensor electroquímico o GC \\
          & $N_{2}O$ emitido* & kg/día & \textit{Inferido vía PGNN (relación N residual)} \\
Productos & Biomasa larval & kg MS/día & Cosecha + secado \\
          & Frass (residuo) & kg MS/día & Balance de masa (Residuo post-proceso) \\
\midrule
\multicolumn{4}{l}{\textit{\textbf{Entradas Indirectas}}} \\
Energía & Electricidad & kWh/día & Medidor de consumo \\
Agua & Agua de humectación & L/día & Contador de flujo \\
\bottomrule
\end{tabularx}
\end{table}

\section{Análisis de Puntos Críticos}

\subsection{El Desafío del $N_2O$ (Óxido Nitroso)}

\begin{center}
\begin{tikzpicture}[node distance=0.5cm]
\node[draw, rectangle, fill=red!10, minimum width=10cm, align=left, rounded corners] (n2o) {
\textbf{¿Qué representa?} \\
Óxido nitroso liberado por desnitrificación incompleta.
};
\node[below=of n2o, draw, rectangle, fill=yellow!10, minimum width=10cm, align=left, rounded corners] (procesos) {
\textbf{Mecanismo:} \\
$NO_3^- \rightarrow NO_2^- \rightarrow \mathbf{N_2O} \rightarrow N_2$ \\
Ocurre en zonas anóxicas dentro del sustrato.
};
\node[below=of procesos, draw, rectangle, fill=orange!10, minimum width=10cm, align=left, rounded corners] (impacto_n2o) {
\textbf{Impacto ambiental:} \\
\textcolor{importante}{\textbf{GWP = 265}} (265 veces más potente que el $CO_2$). \\
Una fuga mínima de este gas puede arruinar la huella de carbono de todo el lote.
};
\end{tikzpicture}
\end{center}

\begin{center}
\begin{tikzpicture}[
    % He eliminado 'drop shadow' para evitar el error de compilación
    larva/.style={draw=none, fill=larva_color, rounded corners=2pt},
    gas_arrow/.style={->, >=latex, thick, shorten >=2pt, shorten <=2pt},
    bad_gas_arrow/.style={->, >=latex, thick, decorate, decoration={snake, amplitude=1.5pt, segment length=5pt}, shorten >=2pt, shorten <=2pt}
]

% --- 1. El Contenedor (Bandeja) con perspectiva ---
\fill[gray!30] (0,0.5) -- (8,0.5) -- (8,5) -- (0,5) -- cycle;
\fill[gray!50] (0,0) -- (8,0) -- (8,0.5) -- (0,0.5) -- cycle;
\fill[gray!40] (8,0) -- (8.5,0.8) -- (8.5,5.3) -- (8,5) -- cycle;
\draw[thick, gray!70] (0,0) -- (8,0) -- (8.5,0.8) -- (8.5,5.3) -- (0.5,5.3) -- (0,5) -- cycle;
\draw[thick, gray!70] (8,0) -- (8,5);

\node[anchor=south west, gray!80, font=\sffamily\bfseries] at (0.2, 5.4) {Bandeja de Bioconversión};

% --- 2. El Sustrato ---

% ZONA ANÓXICA (Inferior)
\begin{scope}
    \clip (0.3, 0.3) to[out=10, in=170] (8.2, 0.6) -- (8.2, 2.5) to[out=190, in=350] (0.3, 2.2) -- cycle;
    \fill[sustrato_anoxico!90!black] (0,0) rectangle (9,6);
    % Textura compacta
    \foreach \i in {1,...,200} \fill[black!40, opacity=0.5] (rnd*8.5, rnd*2.5) circle (0.03);
\end{scope}

% ZONA AERÓBICA (Superior)
\begin{scope}
    \clip (0.3, 2.2) to[out=10, in=170] (8.2, 2.5) -- (8.2, 4.5) to[out=180, in=0] (4, 4.8) to[out=180, in=20] (0.3, 4.2) -- cycle;
    \fill[sustrato_aerobio] (0,0) rectangle (9,6);
    % Textura porosa
    \foreach \i in {1,...,150} \fill[white!60, opacity=0.7] (rnd*8+0.2, rnd*2.5+2.2) circle (rnd*0.08+0.02);
    % Comida
    \foreach \i in {1,...,50} \fill[brown!70!black, opacity=0.6, rounded corners] (rnd*8+0.2, rnd*2.5+2.2) rectangle ++(rnd*0.2, rnd*0.1);
\end{scope}

% --- 3. Larvas ---
\foreach \x/\y/\rot in {1/3.5/20, 2.5/4.1/-10, 4/3.8/5, 5.5/4.2/15, 7/3.6/-20, 3/3/0, 6/2.9/-5} {
    \begin{scope}[shift={(\x,\y)}, rotate=\rot, scale=0.8]
        \fill[larva] (0,0) ellipse (0.4 and 0.15);
        \draw[brown!50, thin] (-0.2, -0.1) -- (-0.2, 0.1);
        \draw[brown!50, thin] (0, -0.13) -- (0, 0.13);
        \draw[brown!50, thin] (0.2, -0.1) -- (0.2, 0.1);
    \end{scope}
}
\node[font=\footnotesize, larva_color!50!black, align=center] at (7.2, 4.3) {\textbf{Larvas}\\\textbf{Activas}};

% --- 4. Etiquetas ---

% Etiqueta Aeróbica
\node[fill=white, draw=sustrato_aerobio, thick, rounded corners, align=left, font=\sffamily\small] (aerobic_label) at (2.5, 5.8) {
    \textbf{ZONA AERÓBICA (Superior)}\\
    \textcolor{gas_bueno}{$\bullet$ Presencia de $O_2$}\\
    $\bullet$ Alta actividad larval
};
\draw[->, thick, sustrato_aerobio] (aerobic_label.south) -- (2.5, 4.5);

% Etiqueta Anóxica
\node[fill=white, draw=sustrato_anoxico, thick, rounded corners, align=left, font=\sffamily\small, anchor=north west] (anoxic_label) at (9, 2.5) {
    \textbf{ZONA ANÓXICA (Inferior)}\\
    \textcolor{gas_malo_n2o}{$\bullet$ Ausencia de $O_2$}\\
    $\bullet$ Sin larvas\\
    $\bullet$ Genera $CH_4$ y $N_2O$
};
\draw[->, thick, sustrato_anoxico] (anoxic_label.west) -- (8.2, 1.5);

% --- 5. Flujos ---
\draw[gas_arrow, gas_bueno, line width=2pt] (1, 6.5) -- (1, 4.8) node[midway, left, font=\bfseries] {$O_2$};
\draw[gas_arrow, gas_neutro, line width=1.5pt] (3.5, 4.5) -- (3.5, 6.5) node[midway, right, font=\bfseries] {$CO_2$};
\draw[bad_gas_arrow, gas_malo_ch4, line width=1.5pt] (5.5, 1.5) -- (5.5, 6.5) node[at end, above, font=\bfseries] {$CH_4$};
\draw[bad_gas_arrow, gas_malo_n2o, line width=1.5pt] (7, 1.0) -- (7, 6.5) node[at end, above, font=\bfseries] {$N_2O$};

\draw[dashed, white, opacity=0.5, thick] (0.3, 2.35) to[out=10, in=170] (8.2, 2.65);

\end{tikzpicture}
\captionof{figure}{Estratificación en la cama de cría. Note las burbujas de metano y óxido nitroso (flechas onduladas) emergiendo de la zona compactada.}
\label{fig:zonas_anoxicas_final}
\end{center}
\subsection{Interpretación Integral del Balance de GEI}

La ecuación para el cálculo del Potencial de Calentamiento Global (GWP) neto es:

\begin{equation}
GWP_{neto} = \underbrace{(M_{CO_2} \cdot 1 + M_{CH_4} \cdot 28 + M_{N_2O} \cdot 265)}_{\text{Emisiones Directas}} + \underbrace{(E_{elec} \cdot FE_{red})}_{\text{Indirectas}} - \underbrace{(M_{larva} \cdot C_{fijado})}_{\text{Crédito}}
\end{equation}

\vspace{0.5cm}

\begin{center}
\begin{tikzpicture}[node distance=1.5cm]
% Nodos circulares
\node[draw, circle, fill=red!20, minimum size=2cm, align=center] (ch4) {\textbf{$CH_4$}\\ \small (GWP 28)};
\node[right=2cm of ch4, draw, circle, fill=red!40, minimum size=2.5cm, align=center] (n2o) {\textbf{$N_2O$}\\ \small (GWP 265)};
\node[right=2cm of n2o, draw, circle, fill=gray!20, minimum size=1.8cm, align=center] (co2) {\textbf{$CO_2$}\\ \small (GWP 1)};

% Nodos de control
\node[below=1.5cm of $(ch4)!0.5!(n2o)$, draw, rectangle, fill=blue!10, rounded corners, minimum width=3cm] (humedad) {\textbf{Humedad > 70\%}};
\node[below=1.5cm of $(n2o)!0.5!(co2)$, draw, rectangle, fill=orange!10, rounded corners, minimum width=3cm] (temp) {\textbf{Temperatura}};

% Flechas
\draw[->, thick, red] (humedad) -- (ch4) node[midway, left, black] {\scriptsize Dispara};
\draw[->, thick, red] (humedad) -- (n2o) node[midway, right, black] {\footnotesize Favorece};
\draw[->, thick, dashed] (temp) -- (n2o);
\draw[->, thick] (temp) -- (co2) node[midway, right] {\footnotesize Acelera};
\end{tikzpicture}
\end{center}

\section{Explicación Detallada del Marco de Evaluación ACV Simplificado}

\subsection{¿Qué es el Análisis de Ciclo de Vida (ACV)?}

El \textbf{Análisis de Ciclo de Vida (ACV)}, también conocido como \textit{Life Cycle Assessment (LCA)} en inglés, es una metodología sistemática para evaluar los impactos ambientales asociados con todas las etapas de la vida de un producto, proceso o servicio. En el contexto de la bioconversión larval, el ACV nos permite cuantificar y comparar la huella ambiental del proceso frente a alternativas tradicionales de manejo de residuos.

\begin{center}
\begin{tikzpicture}[node distance=0.5cm]
\node[draw, rectangle, fill=blue!10, rounded corners, minimum width=12cm, minimum height=1.5cm, align=center, text width=11cm] (def) {
\textbf{ACV Simplificado para Bioconversión} \\
Evalúa los flujos de materia y energía específicamente relacionados con las emisiones de gases de efecto invernadero (GEI) dentro de los límites de la planta (Gate-to-Gate).
};
\end{tikzpicture}
\end{center}

\subsection{Objetivo del Marco de Evaluación}

El objetivo principal de este marco es \textbf{cuantificar las emisiones netas de gases de efecto invernadero} generadas durante el proceso de bioconversión de residuos orgánicos mediante larvas de \textit{Hermetia illucens}, permitiendo:

\begin{itemize}
    \item \textcolor{concepto}{\textbf{Comparación objetiva}} con procesos alternativos (compostaje, digestión anaeróbica).
    \item \textcolor{concepto}{\textbf{Identificación de puntos críticos}} ("Hotspots") donde se generan mayores emisiones.
    \item \textcolor{concepto}{\textbf{Optimización del proceso}} mediante control de variables operativas (aireación, humedad).
    \item \textcolor{concepto}{\textbf{Certificación ambiental}} alineada con estándares (GHG Protocol, ISO 14067).
\end{itemize}

\section{Explicación Detallada de la Tabla de Componentes}

La tabla organizada en tres grandes categorías refleja el \textbf{balance de masa y energía} del sistema:

\begin{table}[H]
\centering
\small
\caption{Componentes del sistema y flujos a cuantificar}
\label{tab:componentes_sistema}
\begin{tabularx}{\textwidth}{@{} l l l X @{}}
\toprule
\textbf{Componente} & \textbf{Flujo} & \textbf{Unidad} & \textbf{Método de cálculo} \\
\midrule
\multicolumn{4}{l}{\textit{\textbf{Entradas}}} \\
Materia Prima & Residuo orgánico & kg MS/día & Balanza + determinación MS (\SI{105}{\celsius}, \SI{24}{\hour}) \\
\midrule
\multicolumn{4}{l}{\textit{\textbf{Salidas Directas}}} \\
Emisiones & $CO_{2}$ emitido & kg/día & Sensor NDIR + integración temporal \\
          & $CH_{4}$ emitido & kg/día & Sensor electroquímico o GC \\
          & $N_{2}O$ emitido* & kg/día & \textit{Inferido vía PGNN (relación N residual)} \\
Productos & Biomasa larval & kg MS/día & Cosecha + secado \\
          & Frass (residuo) & kg MS/día & Balance de masa (Residuo post-proceso) \\
\midrule
\multicolumn{4}{l}{\textit{\textbf{Entradas Indirectas}}} \\
Energía & Electricidad & kWh/día & Medidor de consumo \\
Agua & Agua de humectación & L/día & Contador de flujo \\
\bottomrule
\end{tabularx}
\end{table}

\section{Análisis de Puntos Críticos}

\subsection{El Desafío del $N_2O$ (Óxido Nitroso)}

\begin{center}
\begin{tikzpicture}[node distance=0.5cm]
\node[draw, rectangle, fill=red!10, minimum width=10cm, align=left, rounded corners] (n2o) {
\textbf{¿Qué representa?} \\
Óxido nitroso liberado por desnitrificación incompleta.
};
\node[below=of n2o, draw, rectangle, fill=yellow!10, minimum width=10cm, align=left, rounded corners] (procesos) {
\textbf{Mecanismo:} \\
$NO_3^- \rightarrow NO_2^- \rightarrow \mathbf{N_2O} \rightarrow N_2$ \\
Ocurre en zonas anóxicas dentro del sustrato.
};
\node[below=of procesos, draw, rectangle, fill=orange!10, minimum width=10cm, align=left, rounded corners] (impacto_n2o) {
\textbf{Impacto ambiental:} \\
\textcolor{importante}{\textbf{GWP = 265}} (265 veces más potente que el $CO_2$). \\
Una fuga mínima de este gas puede arruinar la huella de carbono de todo el lote.
};
\end{tikzpicture}
\end{center}

\subsection{Interpretación Integral del Balance de GEI}

La ecuación para el cálculo del Potencial de Calentamiento Global (GWP) neto es:

\begin{equation}
GWP_{neto} = \underbrace{(M_{CO_2} \cdot 1 + M_{CH_4} \cdot 28 + M_{N_2O} \cdot 265)}_{\text{Emisiones Directas}} + \underbrace{(E_{elec} \cdot FE_{red})}_{\text{Indirectas}} - \underbrace{(M_{larva} \cdot C_{fijado})}_{\text{Crédito}}
\end{equation}

\vspace{0.5cm}

\begin{center}
\begin{tikzpicture}[node distance=1.5cm]
% Nodos circulares
\node[draw, circle, fill=red!20, minimum size=2cm, align=center] (ch4) {\textbf{$CH_4$}\\ \small (GWP 28)};
\node[right=2cm of ch4, draw, circle, fill=red!40, minimum size=2.5cm, align=center] (n2o) {\textbf{$N_2O$}\\ \small (GWP 265)};
\node[right=2cm of n2o, draw, circle, fill=gray!20, minimum size=1.8cm, align=center] (co2) {\textbf{$CO_2$}\\ \small (GWP 1)};

% Nodos de control
\node[below=1.5cm of $(ch4)!0.5!(n2o)$, draw, rectangle, fill=blue!10, rounded corners, minimum width=3cm] (humedad) {\textbf{Humedad > 70\%}};
\node[below=1.5cm of $(n2o)!0.5!(co2)$, draw, rectangle, fill=orange!10, rounded corners, minimum width=3cm] (temp) {\textbf{Temperatura}};

% Flechas
\draw[->, thick, red] (humedad) -- (ch4) node[midway, left, black] {\scriptsize Dispara};
\draw[->, thick, red] (humedad) -- (n2o) node[midway, right, black] {\footnotesize Favorece};
\draw[->, thick, dashed] (temp) -- (n2o);
\draw[->, thick] (temp) -- (co2) node[midway, right] {\footnotesize Acelera};
\end{tikzpicture}
\end{center}

\subsection{KPI ambiental propuesto}

\textbf{Ecuacion principal:}

\begin{equation}
\text{GEI}_{neto} \left(\frac{\text{kg CO}_2\text{eq}}{\text{kg residuo}}\right) = \frac{\sum(\text{GEI}_{directos}) + \text{GEI}_{indirectos}}{M_{residuo\_tratado}}
\end{equation}

Donde:
\begin{align*}
\text{GEI}_{directos} &= (M_{\text{CO}_2} \times 1) + (M_{\text{CH}_4} \times 28) + (M_{\text{N}_2\text{O}} \times 265) \\
\text{GEI}_{indirectos} &= (E_{\text{electrica}} \times 0.45 \frac{\text{kg CO}_2\text{eq}}{\text{kWh}}) + \text{transporte}
\end{align*}

\textbf{Metrica de comparacion:}

\begin{equation}
\Delta\text{GEI} = \text{GEI}_{bioconversion} - \text{GEI}_{proceso\_alternativo}
\end{equation}

Procesos alternativos a comparar:
\begin{itemize}
\item Compostaje aerobico: $\sim$0.15--0.30 kg CO\textsubscript{2}eq/kg residuo \cite{ref64}
\item Digestion anaerobica: $\sim$0.05--0.15 kg CO\textsubscript{2}eq/kg residuo \cite{ref54}
\item Incineracion: $\sim$0.80--1.20 kg CO\textsubscript{2}eq/kg residuo (solo emisiones directas)
\end{itemize}

\subsection{Limites del sistema (System Boundary)}

\begin{figure}[H]
\centering
\begin{tikzpicture}[node distance=0.8cm]
\node[draw, rectangle, minimum width=11cm, minimum height=7cm, fill=blue!10] (system) {};

\node[anchor=north west, xshift=0.5cm, yshift=-0.5cm] at (system.north west) (entradas) {\textbf{ENTRADAS:}};
\node[anchor=north west, xshift=0.5cm] at (entradas.south west) (e1) {$\bullet$ Residuo organico (frutas, vegetales, etc.)};
\node[anchor=north west, xshift=0.5cm] at (e1.south west) (e2) {$\bullet$ Larvas inoculo (biomasa inicial)};
\node[anchor=north west, xshift=0.5cm] at (e2.south west) (e3) {$\bullet$ Energia para control ambiental};

\node[anchor=north west, xshift=0.5cm, yshift=-1.5cm] at (e3.south west) (proceso) {\textbf{PROCESO:}};
\node[anchor=north west, xshift=0.5cm] at (proceso.south west) (p1) {$\bullet$ Reactor de bioconversion (batch/continuo)};
\node[anchor=north west, xshift=0.5cm] at (p1.south west) (p2) {$\bullet$ Monitoreo multi-sensor};

\node[anchor=north west, xshift=0.5cm, yshift=-1.5cm] at (p2.south west) (salidas) {\textbf{SALIDAS:}};
\node[anchor=north west, xshift=0.5cm] at (salidas.south west) (s1) {$\bullet$ CO\textsubscript{2}, CH\textsubscript{4}, N\textsubscript{2}O (emisiones)};
\node[anchor=north west, xshift=0.5cm] at (s1.south west) (s2) {$\bullet$ Biomasa larval (producto)};
\node[anchor=north west, xshift=0.5cm] at (s2.south west) (s3) {$\bullet$ Frass (subproducto)};
\node[anchor=north west, xshift=0.5cm] at (s3.south west) (s4) {$\bullet$ Efluente liquido (si aplica)};

\node[anchor=north] at (system.north) {\textbf{SISTEMA DE BIOCONVERSION (enfoque "gate-to-gate")}};
\end{tikzpicture}
\caption{Limites del sistema de bioconversion}
\label{fig:system_boundary}
\end{figure}

\noindent\textbf{Recomendacion critica:} Incluye el \textbf{carbono fijado en biomasa larval} como \textit{secuestro temporal} (factor de -1 en balance), ya que representa carbono que no se mineraliza inmediatamente \cite{ref56}.

\newpage

\section{Fase 2: Mapeo de Variables y Relaciones Biofisicas -- Detalle Mecanistico}

\subsection{Variables medidas vs. latentes}

\begin{table}[H]
\centering
\small % Ayuda a que quepa más contenido cómodamente
\caption{Variables del sistema y sus relaciones con GEI}
\label{tab:variables_gei}
\begin{tabularx}{\textwidth}{@{} l l l X @{}}
\toprule
\textbf{Variable} & \textbf{Sensor / Método} & \textbf{Frecuencia} & \textbf{Relación con GEI / Impacto} \\
\midrule

% GRUPO 1: VARIABLES DIRECTAS
\multicolumn{4}{l}{\textit{\textbf{Variables Directas (Medición física)}}} \\
$CO_2$ & NDIR (0--10k ppm) & 1 min & $\propto$ metabolismo aeróbico + actividad microbiana \\
$CH_4$ & MOS / Electroquímico & 5 min & $\propto$ zonas anaeróbicas (si humedad $>\SI{70}{\percent}$ y hay compactación) \\
Temperatura & Termopar PT100 & 1 min & Controla tasa metabólica ($Q_{10}$): $R(T) = R_{25} \times 2^{\frac{T-25}{10}}$ \\
Humedad Rel. & Capacitivo & 1 min & $>\SI{75}{\percent} \rightarrow$ riesgo anaerobiosis $\rightarrow \uparrow CH_4$ \\
pH & Electrodo combinado & 10 min & $<5.5 \rightarrow$ inhibición; $>8.0 \rightarrow \uparrow NH_3$ (volatilización) \\
Conductividad & Sonda 4 anillos & 10 min & $\propto$ acúmulo de sales ($NH_4^+$, $K^+$) $\rightarrow$ estrés larval \\
$O_2$ & Sensor óptico & 1 min & $<\SI{15}{\percent} \rightarrow$ inicio de condiciones anaeróbicas \\

\addlinespace % Espacio separador elegante

% GRUPO 2: VARIABLES LATENTES
\multicolumn{4}{l}{\textit{\textbf{Variables Latentes (Inferidas vía PGNN)}}} \\
Biomasa & Inferencia (Soft-sensor) & 24 h & $\propto$ Eficiencia de Conversión de Carbono (FCE) \\
$N_2O$ & Inferencia (Soft-sensor) & 1 h & Función compleja $f(pH, NO_3^-\text{res}, H, T)$ \\
Actividad Mic. & Inferencia (Soft-sensor) & 1 h & $\propto$ tasa de mineralización del Carbono \\
\bottomrule
\end{tabularx}
\end{table}
\subsection{Relaciones bioquimicas clave para la PGNN}

\subsubsection{Balance de carbono diferencial}

\begin{equation}
\frac{dC_{total}}{dt} = C_{residuo\_inicial} \times (1 - k_{deg}) - C_{biomasa\_larval} \times \text{FCE} - C_{\text{CO}_2} \times Y_{\text{CO}_2} - C_{\text{CH}_4} \times Y_{\text{CH}_4} - C_{lixiviado}
\end{equation}

Donde:
\begin{itemize}
\item FCE $\approx$ 0.15--0.25 (eficiencia de conversion de carbono)
\item $Y_{\text{CO}_2} \approx$ 0.4--0.6 (rendimiento de CO\textsubscript{2})
\item $Y_{\text{CH}_4} \approx$ 0.01--0.05 en condiciones sub-optimas
\end{itemize}

% AQUI SE APLICA LA CORRECCIÓN DE HYPERREF
\subsubsection{\texorpdfstring{Cinetica de produccion de CH\textsubscript{4} (modelo simplificado)}{Cinetica de produccion de CH4 (modelo simplificado)}}

\begin{equation}
\frac{d[\text{CH}_4]}{dt} = k_{max} \times f(T) \times f(pH) \times f(humedad) \times [\text{Sustrato}_{anoxico}]
\end{equation}

Donde:
\begin{align*}
f(T) &= \exp\left[-\frac{(T - T_{opt})^2}{2\sigma_T^2}\right] && \text{Campana gaussiana ($T_{opt} \approx 35^\circ\text{C}$)} \\
f(pH) &= \frac{1}{1 + \exp[-k_{pH}(pH - pH_{opt})]} && \text{Sigmoide ($pH_{opt} \approx 6.5$--$7.5$)} \\
f(humedad) &= \frac{1}{1 + \exp[-k_h(\text{HR} - \text{HR}_{crit})]} && \text{HR}_{crit} \approx 75\% \text{ para inicio anaerobiosis}
\end{align*}

\subsubsection{Relacion N\textsubscript{2}O-NH\textsubscript{3}-pH (critica para N\textsubscript{2}O no medido)}

\begin{equation}
[\text{N}_2\text{O}] \approx k_{nit} \times [\text{NH}_4^+]_{residual} \times f(pH) \times f(\text{O}_2)
\end{equation}

Con:
\begin{align*}
f(T) &= \exp\left[-\frac{(T - T_{opt})^2}{2\sigma_T^2}\right] && \text{Campana gaussiana } (T_{opt} \approx 35^\circ\text{C}) \\
f(pH) &= \frac{1}{1 + \exp[-k_{pH}(pH - pH_{opt})]} && \text{Sigmoide } (pH_{opt} \approx 6.5\text{--}7.5) \\
f(\text{humedad}) &= \frac{1}{1 + \exp[-k_h(\text{HR} - \text{HR}_{crit})]} && \text{HR}_{crit} \approx 75\% \text{ para inicio anaerobiosis}
\end{align*}

\subsubsection{Efecto de densidad larval en emisiones}

\begin{equation}
\text{Emision}_{\text{GEI}_{total}} = \text{Emision}_{microbiana} \times (1 - \alpha \times \text{Densidad}_{larval})
\end{equation}

Donde $\alpha \approx$ 0.2--0.4: las larvas consumen sustrato antes de que microbios lo mineralicen $\rightarrow$ Menor tiempo de residencia del carbono = menor mineralizacion a GEI.

\noindent\textbf{Hallazgo clave de literatura:} Bioconversion con larvas reduce emisiones de GEI vs. compostaje en 30--60\% debido a menor tiempo de residencia y consumo directo del carbono por larvas \cite{ref56,ref64}.

\newpage

\section{Fase 3: Arquitectura de la PGNN -- Detalle de Implementacion}

\subsection{Estructura hibrida propuesta}

\begin{figure}[H]
\centering
% Se usa >=stealth, lo cual requiere que tikz maneje bien el caracter >.
% Con \usetikzlibrary{babel} cargado, esto ya no genera error.
\begin{tikzpicture}[node distance=0.8cm, >=stealth]
\node[draw, rectangle, minimum width=9cm, minimum height=0.8cm, fill=green!20, align=center] (input) {ENTRADAS (t): [CO\textsubscript{2}, CH\textsubscript{4}, T, HR, pH, CE, O\textsubscript{2}, t]};

\node[draw, rectangle, below=of input, minimum width=9cm, minimum height=1.2cm, fill=blue!20, align=center] (physics) {MODULO FISICO (ODE solver)};

\node[draw, rectangle, below=of physics, minimum width=9cm, minimum height=0.8cm, fill=yellow!20, align=center] (residual) {RESIDUO = Datos\textsubscript{reales} - Pred\textsubscript{fisica}};

\node[draw, rectangle, below=of residual, minimum width=9cm, minimum height=1.5cm, fill=red!20, align=center] (nn) {RED NEURONAL (LSTM o TCN)};

\node[draw, rectangle, below=of nn, minimum width=9cm, minimum height=0.8cm, fill=purple!20, align=center] (output) {SALIDA: GEI\textsubscript{total}(t) + incertidumbre $\sigma$(t)};

\draw[->, thick] (input) -- (physics);
\draw[->, thick] (physics) -- (residual);
\draw[->, thick] (residual) -- (nn);
\draw[->, thick] (nn) -- (output);
\end{tikzpicture}
\caption{Arquitectura hibrida PGNN para evaluacion ambiental}
\label{fig:pgnn_architecture}
\end{figure}

\subsection{Formulacion matematica de la perdida hibrida}

\begin{equation}
\text{Loss}_{total} = \alpha \cdot \text{MSE}_{datos} + \beta \cdot \text{MSE}_{fisico} + \gamma \cdot \text{MSE}_{conservacion}
\end{equation}

Donde:

\begin{align*}
\text{MSE}_{datos} &= \frac{1}{N} \sum \|y_{pred} - y_{medido}\|^2 && \text{Ajuste a datos reales} \\
\text{MSE}_{fisico} &= \frac{1}{N} \sum \left\|\frac{dC}{dt}_{pred} - f_{fisica}(x)\right\|^2 && \text{Penalizacion por violar ODEs} \\
\text{MSE}_{conservacion} &= \lambda \cdot |C_{entrada} - (C_{biomasa} + C_{\text{CO}_2} + C_{\text{CH}_4} + C_{residual})|^2 && \text{Balance de masa global ($\approx$ 0)}
\end{align*}

\textbf{Valores sugeridos para hiperparametros:}
\begin{itemize}
\item $\alpha = 1.0$ (peso datos)
\item $\beta = 0.3$--$0.5$ (peso fisica -- ajustar segun calidad de modelo mecanistico)
\item $\gamma = 0.2$ (peso conservacion)
\item $\lambda = 10$--$100$ (factor de penalizacion fuerte para balance de masa)
\end{itemize}

\subsection{Arquitectura neuronal recomendada}

\begin{table}[H]
\centering
\caption{Arquitectura neuronal para PGNN}
\label{tab:nn_architecture}
\begin{tabular}{|l|l|c|l|c|l|}
\hline
\textbf{Capa} & \textbf{Tipo} & \textbf{Unidades} & \textbf{Activacion} & \textbf{Dropout} & \textbf{Notas} \\ \hline
Input & - & 9 & - & - & Normalizar a $\mu=0$, $\sigma=1$ \\ \hline
LSTM 1 & Recurrente & 64 & tanh & 0.2 & Captura memoria temporal \\ \hline
LSTM 2 & Recurrente & 32 & tanh & 0.2 & - \\ \hline
Dense 1 & Feedforward & 32 & ReLU & 0.3 & - \\ \hline
Dense 2 (salida) & Feedforward & 3 & lineal & - & [CO\textsubscript{2}, CH\textsubscript{4}, N\textsubscript{2}O] \\ \hline
Dense 3 (incertidumbre) & Feedforward & 3 & softplus & - & $\sigma_{\text{CO}_2}$, $\sigma_{\text{CH}_4}$, $\sigma_{\text{N}_2\text{O}}$ \\ \hline
\end{tabular}
\end{table}

\noindent\textbf{Alternativa si datos escasos ($<$1000 puntos):} Usar \textbf{Neural ODE} (Chen et al. 2018) donde la red aprende directamente los parametros de las ODEs \cite{ref4}.

\subsection{Protocolo de entrenamiento}

\begin{lstlisting}[caption=Pseudocodigo PyTorch para entrenamiento PGNN]
# Pseudocodigo PyTorch
for epoch in range(epochs):
    # Paso 1: Forward pass con datos reales
    y_pred, sigma = pgnn_model(X_batch)

    # Paso 2: Calcular perdida de datos (con incertidumbre)
    loss_data = torch.mean(0.5 * torch.log(sigma**2) +
                          0.5 * ((y_true - y_pred)**2) / sigma**2)

    # Paso 3: Forward pass con condiciones iniciales para ODE
    y_phys = physics_baseline(X_batch)  # Resolver ODE con parametros actuales

    # Paso 4: Calcular perdida fisica
    loss_phys = MSE(y_pred - y_phys_residual, NN_residual)

    # Paso 5: Balance de masa global
    loss_mass = mass_balance_violation(X_batch, y_pred)

    # Paso 6: Backpropagation
    total_loss = alpha*loss_data + beta*loss_phys + gamma*loss_mass
    total_loss.backward()
    optimizer.step()
\end{lstlisting}

\newpage

\section{Fase 4: Validacion y Comparacion Ambiental -- Detalle Metodologico}

\subsection{Diseno experimental para recoleccion de datos}

\begin{table}[H]
\centering
\caption{Diseno experimental para validacion}
\label{tab:experimental_design}
\begin{tabular}{|l|c|l|l|}
\hline
\textbf{Condicion experimental} & \textbf{Replicas} & \textbf{Variables controladas} & \textbf{Objetivo} \\ \hline
\textbf{Densidad larval baja} & 3 & 5 larvas/g residuo & Establecer baseline microbiano \\ \hline
\textbf{Densidad larval media} & 3 & 15 larvas/g residuo & Condicion optima operacional \\ \hline
\textbf{Densidad larval alta} & 3 & 30 larvas/g residuo & Evaluar estres y competencia \\ \hline
\textbf{Sin larvas (control)} & 3 & 0 larvas/g & Cuantificar emisiones puramente microbianas \\ \hline
\textbf{Variacion de sustrato} & 2$\times$3 & Fruta vs. vegetal & Evaluar efecto composicion C/N \\ \hline
\end{tabular}
\end{table}

\textbf{Duracion minima por replica:} 14 dias (ciclo completo de desarrollo larval)

\textbf{Frecuencia de muestreo:}
\begin{itemize}
\item Gases: continuo (1 min)
\item Parametros fisico-quimicos: cada 4h
\item Muestreo destructivo (biomasa larval, analisis residual): dias 0, 3, 7, 10, 14
\end{itemize}

\subsection{Validacion cruzada especifica para PGNN}

\begin{lstlisting}[caption=Estrategia de validacion cruzada]
Estrategia propuesta:

Fold 1: Entrenar con densidades baja + media -> Validar en alta
Fold 2: Entrenar con fruta -> Validar en vegetal
Fold 3: Entrenar con T=25C -> Validar en T=35C
Fold 4: Entrenar con dias 0-7 -> Validar en dias 8-14 (extrapolacion temporal)

Metricas de exito:
- R^2 > 0.85 en validacion (datos no vistos)
- Error en balance de masa < 5%
- Incertidumbre calibrada: 95% de puntos reales dentro de +-2sigma_pred
\end{lstlisting}

\subsection{Comparacion con procesos alternativos (protocolo riguroso)}

\textbf{Paso 1:} Calcular GEI\textsubscript{neto} para bioconversion (tu sistema)

\textbf{Paso 2:} Obtener GEI para alternativas:
\begin{itemize}
\item \textbf{Opcion A (literatura):} Usar valores de meta-analisis \cite{ref54,ref64} con intervalos de confianza
\item \textbf{Opcion B (experimental):} Ejecutar compostaje paralelo en reactor controlado con mismos residuos
\end{itemize}

\textbf{Paso 3:} Analisis estadistico:

\begin{lstlisting}[caption=Prueba de significancia estadistica]
# Prueba de significancia (si tienes replicas)
from scipy import stats

bioconversion_gei = [0.12, 0.14, 0.11, 0.13]  # kg CO2eq/kg
compostaje_gei    = [0.28, 0.31, 0.26, 0.29]

t_stat, p_value = stats.ttest_ind(bioconversion_gei, compostaje_gei)
# p < 0.05 -> diferencia significativa
\end{lstlisting}

\textbf{Paso 4:} Analisis de sensibilidad global (Sobol indices):

\begin{lstlisting}[caption=Analisis de sensibilidad global]
# Identificar que variables operativas impactan mas el GEI_neto
S_T_humedad   = 0.35  # Contribucion total de humedad a varianza de GEI
S_T_densidad  = 0.28
S_T_Temp      = 0.22
S_T_pH        = 0.15

-> Conclusion: Controlar humedad es critico para minimizar GEI
\end{lstlisting}

\newpage

\section{Fase 5: Implementacion Practica -- Roadmap Detallado}

\subsection{Cronograma de 6 meses}

\begin{table}[H]
\centering
\caption{Cronograma de implementacion}
\label{tab:cronograma}
\begin{tabular}{|c|p{7.5cm}|p{4.5cm}|}
\hline
\textbf{Mes} & \textbf{Actividades} & \textbf{Entregables} \\ \hline
\textbf{1} & $\bullet$ Diseno reactor + seleccion sensores\\$\bullet$ Adquisicion equipos\\$\bullet$ Protocolo muestreo & Lista de materiales, diagrama P\&ID \\ \hline
\textbf{2} & $\bullet$ Montaje sistema\\$\bullet$ Calibracion sensores (CO\textsubscript{2}, CH\textsubscript{4}, T, pH)\\$\bullet$ Primeras corridas sin larvas & Datos baseline microbianos \\ \hline
\textbf{3} & $\bullet$ Corridas con larvas (3 densidades $\times$ 3 replicas)\\$\bullet$ Almacenamiento datos estructurado (time-series DB) & Dataset completo $\geq$ 5,000 puntos temporales \\ \hline
\textbf{4} & $\bullet$ Desarrollo PGNN (PyTorch/TensorFlow)\\$\bullet$ Entrenamiento + validacion cruzada\\$\bullet$ Analisis de sensibilidad & Modelo .pt/.h5 + reporte metricas \\ \hline
\textbf{5} & $\bullet$ Integracion en dashboard (Grafana/Plotly)\\$\bullet$ Calculo automatico KPI GEI\textsubscript{neto}\\$\bullet$ Comparacion vs. literatura & Interfaz operativa + reporte preliminar \\ \hline
\textbf{6} & $\bullet$ Validacion campo (si aplica)\\$\bullet$ Documentacion para ACV completo\\$\bullet$ Preparacion articulo/patente & Reporte tecnico final + roadmap escalado \\ \hline
\end{tabular}
\end{table}

\subsection{Lista de equipos minimos (presupuesto estimado)}

\begin{table}[H]
\centering
\caption{Lista de equipos minimos}
\label{tab:equipos_minimos}
\begin{tabular}{|l|l|r|}
\hline
\textbf{Equipo} & \textbf{Especificacion minima} & \textbf{Costo aprox. (USD)} \\ \hline
Sensor CO\textsubscript{2} & NDIR, rango 0--10,000 ppm, $\pm$50 ppm & \$150 \\ \hline
Sensor CH\textsubscript{4} & Electoquimico, 0--500 ppm, $\pm$10 ppm & \$200 \\ \hline
Sensor O\textsubscript{2} & Optico, 0--25\%, $\pm$0.1\% & \$180 \\ \hline
Termohigrometro & PT100 + capacitivo, $\pm$0.3$^\circ$C / $\pm$3\% HR & \$120 \\ \hline
pH-metro & Electrodo combinado, rango 2--12 & \$250 \\ \hline
Conductivimetro & 4 anillos, rango 0--20 mS/cm & \$100 \\ \hline
DAQ & Arduino Mega + shield Ethernet o Raspberry Pi 4 & \$80 \\ \hline
Reactor & Contenedor plastico 20L con ventilacion controlada & \$50 \\ \hline
\textbf{TOTAL} & & \textbf{\$1,130} \\ \hline
\end{tabular}
\end{table}

\noindent\textbf{Alternativa low-cost:} Usar sensores MQ-135 (CO\textsubscript{2}) y MQ-4 (CH\textsubscript{4}) con calibracion cruzada contra estandares certificados (error $\pm$15\%, aceptable para prototipo).

\subsection{Integracion con estandares de reporte ambiental}

\begin{table}[H]
\centering
\caption{Integracion con estandares ambientales}
\label{tab:estandares_ambientales}
\begin{tabular}{|l|l|l|}
\hline
\textbf{Estandar} & \textbf{Requisito} & \textbf{Como cumplirlo con tu sensor virtual} \\ \hline
\textbf{GHG Protocol} & Reporte Scope 1 (emisiones directas) & Tu PGNN entrega CO\textsubscript{2}, CH\textsubscript{4}, N\textsubscript{2}O directos \\ \hline
\textbf{ISO 14067} & Huella de carbono de productos & Incluir biomasa larval como producto con carbono fijado \\ \hline
\textbf{PAS 2050} & ACV de bienes y servicios & Tu KPI GEI\textsubscript{neto} es input directo \\ \hline
\textbf{Science Based Targets} & Reduccion absoluta de GEI & Demostrar $\Delta$GEI vs. baseline sectorial \\ \hline
\end{tabular}
\end{table}

\newpage

\section{Acciones Inmediatas (Proximas 72 Horas)}

\begin{enumerate}
\item \textbf{Define tu sustrato especifico} (ej. "residuos de fruteria con C/N=25") $\rightarrow$ determina cineticas esperadas
\item \textbf{Selecciona larva objetivo} (\textit{Hermetia illucens} vs. \textit{Tenebrio molitor}) $\rightarrow$ afecta densidad optima y perfil de emisiones
\item \textbf{Disena matriz experimental} (3 densidades $\times$ 3 replicas $\times$ 14 dias) $\rightarrow$ minimo 126 observaciones temporales
\item \textbf{Prioriza medicion de CH\textsubscript{4}} $\rightarrow$ es la variable critica que diferencia bioconversion eficiente (bajo CH\textsubscript{4}) vs. ineficiente (alto CH\textsubscript{4} por anaerobiosis)
\item \textbf{Implementa balance de carbono simple en Excel} $\rightarrow$ antes de PGNN, verifica que tus mediciones cierren balances (error $<$10\%)
\end{enumerate}

\newpage

\begin{thebibliography}{99}
\bibitem{ref4} Chen, R. T. Q., Rubanova, Y., Bettencourt, J., \& Duvenaud, D. K. (2018). Neural ordinary differential equations. In \textit{Advances in neural information processing systems} (pp. 6571--6583).

\bibitem{ref54} Referencia sobre digestion anaerobica y GEI.

\bibitem{ref56} Referencia sobre reduccion de GEI en bioconversion vs compostaje.

\bibitem{ref59} Referencia sobre relacion N\textsubscript{2}O-NH\textsubscript{3}-pH.

\bibitem{ref64} Referencia sobre GEI en compostaje aerobico.
\end{thebibliography}

\end{document}
